%% LyX 2.4.4 created this file.  For more info, see https://www.lyx.org/.
%% Do not edit unless you really know what you are doing.
\documentclass[english]{article}
\usepackage[T1]{fontenc}
\usepackage[latin9]{inputenc}
\usepackage{units}
\usepackage{enumitem}
\usepackage{amsmath}
\usepackage{amsthm}
\usepackage{graphicx}
\usepackage{geometry}
\geometry{verbose,tmargin=1in,bmargin=1in,lmargin=1in,rmargin=1in}

\makeatletter
%%%%%%%%%%%%%%%%%%%%%%%%%%%%%% Textclass specific LaTeX commands.
\numberwithin{equation}{section}
\numberwithin{figure}{section}
\newlength{\lyxlabelwidth}      % auxiliary length 

%%%%%%%%%%%%%%%%%%%%%%%%%%%%%% User specified LaTeX commands.
\usepackage{fancyhdr}
\usepackage{lastpage}
\pagestyle{fancy}
\usepackage{enumitem}
\usepackage{ifthen}
%sets math fonts to be more readable
\everymath=\expandafter{\the\everymath\displaystyle}
%\DeclareMathSizes{10}{10}{10}{10}
\usepackage{multicol}
% Added by lyx2lyx
\setlength{\parskip}{\smallskipamount}
\setlength{\parindent}{0pt}

\makeatother

\usepackage{babel}
\begin{document}
\renewcommand{\author}{Dr.\,James Ashe}

\newcommand{\classNum}{231}

\newcommand{\classSec}{C}

\newcommand{\nameType}{Name:}

\newcommand{\examType}{Exam 2}

\renewcommand{\date}{November 9, 2012}

%%First page's header%%%%%%%%%%%%%%%%%%%%%%
\fancypagestyle{firststyle} %%header settings for first pagr
{
	\fancyhead[L]{MTH \classNum{}(\classSec{}) \author{}\\
	\textbf{\examType{}} \date{}}
	\fancyhead[C]{\textbf{\nameType}}
	\setlength{\headsep}{14pt}
}
%%%%%%%%%%%%%%%%%%%%%%%%%%%%%%%%%%%%%%%%%%

%%Next page's header%%%%%%%%%%%%%%%%%%%%%%%
\setlist{nolistsep}
\lhead{\classNum{}(\classSec{})} 
\chead{\textbf{\nameType}} 
\cfoot{\thepage{}~of~\pageref{LastPage}} 
\setlength{\headsep}{5pt} 
%%%%%%%%%%%%%%%%%%%%%%%%%%%%%%%%%%%%%%%%%%%

%%Various settings; see the preamble for other settings%%%%%
%%\setenumerate[0]{leftmargin=12pt,itemindent=0pt,labelwidth=0pt}
\renewcommand{\labelenumi}{\textbf{\arabic{enumi}.}}
\renewcommand{\labelenumii}{\textbf{\alph{enumii}.}}
\thispagestyle{firststyle}
%%%%%%%%%%%%%%%%%%%%%%%%%%%%%%%%%%%%%%%%%%

\emph{Find the displacement over the given interval.}
\begin{enumerate}[resume]
\item $v(t)=\cos\left(\frac{\pi}{2}t\right)$ on $[0,3]$; $s(0)=0$.\vfill
\end{enumerate}
\emph{Determine the position function.}
\begin{enumerate}[resume]
\item An object is moving at a velocity of $(t+9)^{\nicefrac{1}{2}}$ with
an initial position of $s(0)=20$ at time $t$.\vfill
\end{enumerate}
\emph{Solve the problem.}
\begin{enumerate}[resume]
\item A projectile is launched vertically from the ground at $t=0$ where
$t$ is seconds, and its velocity in flight (in $\nicefrac{\mbox{m}}{\mbox{s}}$)
is given by $v(t)=20-10t$.

\begin{enumerate}
\item Find the position at time $t$ for $0\leq t\leq4$.
\item Find the total distance traveled from $t=0$ to $t=4$.\vspace{3cm}\vfill\newpage
\end{enumerate}
\end{enumerate}
\emph{Use calculus to find the area of the region described.}
\begin{enumerate}[resume]
\item The region bounded by $y=x$ and $y=x^{2}-4x$.\\
\includegraphics[width=3cm]{D:/home/jim/JamesAshe/JCSU/Calc_2_MTH_232/images/x_and_x2-4x.png}\vfill
\item The region bounded by $y=\frac{1}{2}x^{2}$ and $y=-x^{2}+6$.\\
\includegraphics[width=3cm]{D:/home/jim/JamesAshe/JCSU/Calc_2_MTH_232/images/1/2x2_and-x2+6.png}\vfill
\item The region in the first quadrant bounded by $-x+1$ and $x-1$ from
$0\leq x\leq2$.\\
\includegraphics[width=3cm]{D:/home/jim/JamesAshe/JCSU/Calc_2_MTH_232/images/1-x_and_x-1.png}\vfill\newpage
\end{enumerate}
\emph{Use the general slicing method and calculus to find the volume
of the solid.}
\begin{enumerate}[resume]
\item The rectangular box with dimensions a height of $3$, width of $4$,
and length of $10$.\vfill
\end{enumerate}
\emph{Let $R$ be the region bounded by the following curves. Use
the disk method to find the volume of the solid generated when $R$
is revolved about the $x$-axis.}
\begin{enumerate}[resume]
\item $y=\sqrt{x}$, $x=0$, and $x=9$.\vfill
\item $y=\sqrt{\sin4x}$, $y=0$, $0\leq x\leq\frac{\pi}{4}$.\vfill\newpage
\end{enumerate}
\emph{Find the length of the curve.}
\begin{enumerate}[resume]
\item $y=2x^{\nicefrac{3}{2}}$ from $x=0$ to $x=\frac{5}{4}$.\vfill
\item $y=\left(4-x^{\nicefrac{2}{3}}\right)^{\nicefrac{3}{2}}$ from $x=0$
to $x=8$.\vfill
\end{enumerate}
\emph{BONUS: Use calculus and the general slicing method to find a
general formula for the volume of a pyramid with a square base.}\vfill
\end{document}
